\section*{الفصل \ref{ch:g8:solids} --- الأهرام والمخاريط}

\begin{solution}{exo:g8:solids:1}
\omterm{def:g4:geometry:polygon}{رباعي} الوجوه: $4$ \omterm{def:g5:solids:def}{أوجه}، و $6$ \omterm{def:g5:solids:def}{أحرف}، و $4$ \omterm{def:g5:solids:def}{رؤوس}
($4 + 4 = 6 + 2$). والقاعدة السداسية: $7$ \omterm{def:g5:solids:def}{أوجه}، و $12$ حرفًا، و $7$
\omterm{def:g5:solids:def}{رؤوس} ($7 + 7 = 12 + 2$).
\end{solution}

\begin{solution}{exo:g8:solids:2}
$V = \frac13 \times 36 \times 10 = 120$ سم$^3$.

والقاعدة المستطيلة: $B = 24$ سم$^2$، إذن
$V = \frac13 \times 24 \times 7.5 = 60$ سم$^3$.
\end{solution}

\begin{solution}{exo:g8:solids:3}
$V = \frac13 \pi \times 25 \times 9 = 75\pi \approx 236$ سم$^3$.
\end{solution}

\begin{solution}{exo:g8:solids:4}
\omterm{def:g7:areas:prism}{الأسطوانة}: $\pi \times 16 \times 12 = 192\pi \approx 603$ سم$^3$.
\omterm{def:g8:solids:def}{والمخروط}: $\frac{192\pi}{3} = 64\pi \approx 201$ سم$^3$. والنسبة: \omterm{def:g8:solids:def}{فالمخروط}
ثلث \omterm{def:g7:areas:prism}{الأسطوانة}.
\end{solution}

\begin{solution}{exo:g8:solids:5}
$56 = \frac13 \times B \times 8$، إذن $B = \frac{56 \times 3}{8} = 21$
سم$^2$.
\end{solution}

\begin{solution}{exo:g8:solids:6}
ينفرد السطح الجانبي في قطاع قرص نصف قطره هو
\emph{الطول المائل} $s$ (وهو بُعد القمّة عن الحافة ---
فتلك \omterm{def:g6:lines:objects}{القطعة} تقع على السطح)، لا الارتفاع $h$.
\end{solution}

\begin{solution}{exo:g8:solids:7}
$V = \frac13 \times 35^2 \times 22 = \frac13 \times 1225 \times 22
= \frac{26\,950}{3} \approx 8\,983$ م$^3$ --- أي نحو $9\,000$ متر
\omterm{def:g5:solids:def}{مكعب}.
\end{solution}

\begin{solution}{exo:g8:solids:8}
الارتفاع: $h = \sqrt{13^2 - 5^2} = \sqrt{169 - 25} = \sqrt{144} = 12$
سم. \omterm{def:g6:measure:volume}{والحجم}: $V = \frac13 \pi \times 25 \times 12 = 100\pi$ سم$^3$.
\end{solution}

\begin{solution}{exo:g8:solids:9}
\omterm{def:g3:shapes:shapes}{نصف القطر} نفسه، إذن \omterm{def:g6:measure:volume}{حجم} \omterm{def:g8:solids:def}{المخروط} ثلث \omterm{def:g6:measure:volume}{حجم}
\omterm{def:g7:areas:prism}{الأسطوانة} \emph{من أجل الارتفاع نفسه}: فيبلغ السائل
$\frac{9}{3} = 3$ سم في \omterm{def:g7:areas:prism}{الأسطوانة}.
\end{solution}

\begin{solution}{exo:g8:solids:10}
الرمل $=$ \omterm{def:g8:solids:def}{مخروط} واحد: $V = \frac13 \pi \times 4 \times 4.5 = 6\pi
\approx 18.8$ سم$^3$. والساعة الرملية تسع \omterm{def:g8:solids:def}{مخروطين} كهذا، إذن
يملأ الرمل نصف \omterm{def:g6:measure:volume}{الحجم} الداخلي الكلي بالضبط.
\end{solution}

\begin{solution}{exo:g8:solids:11}
\emph{1.} نعم: فالصيغة $V = \frac13 Bh$ تصحّ من أجل أيّ موضع
للقمّة، ما دام $h$ هو البُعد العمودي عن مستوي
القاعدة. و $V = \frac13 \times 36 \times 8 = 96$ سم$^3$.

\emph{2.} \omterm{def:g4:geometry:circle}{قطر} القاعدة: $\sqrt{6^2 + 6^2} = 6\sqrt2$ سم. وأطول
حرف جانبي هو وتر \omterm{def:g6:shapes:triangles}{مثلث قائم الزاوية} ضلعا قائمته
$6\sqrt2$ (في مستوي القاعدة) و $8$ (عموديًّا):
\[
\sqrt{(6\sqrt2)^2 + 8^2} = \sqrt{72 + 64} = \sqrt{136} \approx 11.7
\text{ سم}.
\]
\end{solution}

\begin{solution}{pb:g8:solids:1}
\textbf{1.} الرأس $G$ ينتمي إلى الوجه الأعلى $EFGH$ وإلى الوجهين
الجانبيين $BCGF$ و $DCGH$. \omterm{def:g5:solids:def}{والأوجه} الثلاثة التي \emph{لا}
تحوي $G$ هي الأسفل $ABCD$، والأمام $ABFE$، والوجه
الأيسر $ADHE$ (وهي كلها تشترك في الرأس $A$، وهو الركن
المقابل للرأس $G$). والأهرام الثلاثة هي $ABCDG$ و $ABFEG$ و
$ADHEG$: لكل منها وجه مربع من \omterm{def:g5:solids:def}{المكعب} قاعدةً والرأس
البعيد $G$ قمّةً. وكل واحد \omterm{def:g8:solids:def}{هرم} ذو قاعدة مربعة: $5$ \omterm{def:g5:solids:def}{أوجه}،
و $8$ \omterm{def:g5:solids:def}{أحرف}، و $5$ \omterm{def:g5:solids:def}{رؤوس} (\cref{ex:g8:solids:count}).

\textbf{2.} \omterm{def:g4:geometry:circle}{القطر} الطويل $[AG]$ يصل الرأسين الوحيدين
اللذين لا ينتميان إلى أيّ من القواعد الثلاث. وإدارة بثلث دورة
حول \omterm{def:g6:lines:objects}{المستقيم} $(AG)$ ترسل \omterm{def:g5:solids:def}{المكعب} إلى نفسه (\omterm{def:g5:solids:def}{فالأحرف} الثلاثة
الخارجة من $A$ --- نحو $B$ و $D$ و $E$ --- تُخلط في
دورة، وكذلك \omterm{def:g5:solids:def}{الأحرف} الثلاثة الواصلة إلى $G$). وهي تحمل
القاعدة $ABCD$ إلى $ADHE$، ثم $ADHE$ إلى $ABFE$، وترجع:
فالقمم ثابتة عند $G$، والقواعد تدور. ومنه يُحمل كل \omterm{def:g8:solids:def}{هرم}
على الذي يليه: فالثلاثة نسخ متطابقة.

\textbf{3.} ثلاث قطع متطابقة تملأ \omterm{def:g5:solids:def}{المكعب} تتقاسم
حجمه بالتساوي:
\[
V_{\text{الهرم}} = \frac{a^3}{3} .
\]

\textbf{4.} \omterm{def:g8:solids:def}{الهرم} $ABCDG$ قاعدته $ABCD$، ومساحتها $a^2$،
وقمّته $G$ تقع عموديًّا فوق الركن $C$، على اِرتفاع
$CG = a$ (وهو حرف عمودي من \omterm{def:g5:solids:def}{المكعب}). وكما في
\cref{exo:g8:solids:11}، تنطبق الصيغة بالارتفاع
العمودي $h = a$:
\[
V = \frac13 \times a^2 \times a = \frac{a^3}{3} ,
\]
في توافق تامّ مع السؤال 3 --- فالتفكيك
والصيغة يؤكد كل منهما الآخر.

\textbf{5.} من أجل $a = 6$:
$V = \frac{6^3}{3} = \frac{216}{3} = 72$ سم$^3$ لكل واحد. و
تجربة سكب الماء هي هذا التفكيك مسيَّلًا:
\omterm{def:g7:areas:prism}{فالموشور} فوق القاعدة $ABCD$ باِرتفاع $a$ هو \omterm{def:g5:solids:def}{المكعب} نفسه،
وهو يسع ثلاث ملآت من \omterm{def:g8:solids:def}{الهرم} بالضبط --- وهنا تتجمّع الملآت
الثلاث هندسيًّا أيضًا في \omterm{def:g5:solids:def}{المكعب}.

\textbf{6.} وصل المركز $O$ بأركان كل وجه
الأربعة ينتج ستة أهرام، واحدًا لكل وجه: قاعدة مربعة، وقمّة $O$.
والمركز موضوع بالتساوي بالنسبة إلى \omterm{def:g5:solids:def}{الأوجه} الستة (فأيّ
دوران أو تماثل \omterm{def:g5:solids:def}{للمكعب} يثبّت $O$ ويبدّل
\omterm{def:g5:solids:def}{الأوجه})، إذن الأهرام الستة متطابقة. واِرتفاع كل واحد هو
بُعد $O$ عن وجه: أي نصف الضلع، $\frac{a}{2}$.

\textbf{7.} ست قطع متطابقة تتقاسم \omterm{def:g5:solids:def}{المكعب}:
أي $V = \frac{a^3}{6}$ لكل واحدة.

\textbf{8.} بالصيغة: \omterm{def:g6:measure:area}{مساحة} القاعدة $a^2$، والارتفاع $\frac a2$:
\[
V = \frac13 \times a^2 \times \frac{a}{2} = \frac{a^3}{6} .
\]
والجوابان متوافقان.

\textbf{9.} يخصّ التفكيكان أهرامًا ذات نسب
خاصّة جدًّا منحوتة من \omterm{def:g5:solids:def}{مكعب}؛ أما \omterm{def:g8:solids:def}{الهرم} النحيل، أو المائل،
أو \omterm{def:g8:solids:def}{المخروط}، فلا يمكن تركيبه بهذه الطريقة (فلا ثلاثة \omterm{def:g8:solids:def}{مخاريط} تملأ
\omterm{def:g7:areas:prism}{أسطوانة}). ولهذا كانت \cref{thm:g8:solids:volume}
\emph{مقبولة} في هذا المستوى: فالتفكيكات وتجربة
سكب الماء تجعل العامل $\frac13$ معقولًا تمامًا،
والبرهان العامّ الأمين --- وهو التكامل --- يُعطى في
مجلّد المرحلة الثانوية.

\textbf{10.} $V = \frac{6^3}{6} = 36$ سم$^3$؛ وبالصيغة،
$V = \frac13 \times 36 \times 3 = 36$ سم$^3$. وهما متوافقان.

\textbf{11.} في المثلث $SAB$، يوصَل بين منتصفَي الضلعين
$[SA]$ و $[SB]$ \omterm{def:g6:lines:objects}{بقطعة} موازية \omterm{def:g6:lines:objects}{للمستقيم} $(AB)$ و
طولها $\frac{AB}{2} = \frac{c}{2}$
(\cref{thm:g8:midpoints:direct}). والشيء نفسه يصحّ في كل وجه
جانبي، إذن تكوّن المنتصفات الأربعة مربعًا طول ضلعه $\frac c2$،
وأضلاعه موازية للقاعدة. وعلى \omterm{def:g6:lines:objects}{المستقيم} العمودي المارّ
بالقمّة، يمرّ المقطع في منتصف الارتفاع (بمبرهنة
المنتصفين من جديد، في مثلث يمرّ بالقمّة $S$ ومركز القاعدة ورأس
من القاعدة): فيقع مستوي \omterm{def:g6:lines:objects}{القطع} على اِرتفاع $\frac h2$.

\textbf{12.} \omterm{def:g6:lines:objects}{القطعة} العليا \omterm{def:g8:solids:def}{هرم} منتظم ذو قاعدة مربعة طول
ضلعها $\frac c2$ واِرتفاعه $\frac h2$ (فقاعدته هي المقطع،
وقمّته هي $S$). وحجمه
\[
\frac13 \times \left(\frac c2\right)^2 \times \frac h2
= \frac13 \times \frac{c^2}{4} \times \frac h2
= \frac{1}{8} \times \frac{c^2 h}{3}
= \frac{V}{8} :
\]
أي ثُمن الأصلي --- لا نصفه.

\textbf{13.} يحتفظ جذع \omterm{def:g8:solids:def}{الهرم} \omterm{def:g3:division:remainder}{بالباقي}:
$V - \frac V8 = \frac{7V}{8}$. فقطع نصف الارتفاع يتقاسم
\omterm{def:g6:measure:volume}{الحجم} بنسبة $1 : 7$ --- فالبلاطة السفلى المتواضعة المظهر
تحمل سبعة أمثال \omterm{def:g6:measure:volume}{حجم} القمّة المدبّبة.

\textbf{14.} يقع تحت نصف الارتفاع $\frac78$ من \omterm{def:g6:measure:volume}{الحجم}. ومع
$V = \frac13 \times 35^2 \times 22 = \frac{26\,950}{3} \approx
8\,983$ م$^3$ (\cref{exo:g8:solids:7}):
\[
\frac78 \times \frac{26\,950}{3} \approx 7\,860 \approx
7\,900 \text{ م}^3
\]
إلى أقرب مئة --- فحملة «النصف السفلي» تنظّف
نحو ثمانية أمثال \omterm{def:g6:measure:volume}{حجم} زجاج الحملة العليا.

\textbf{15.} مضاعفة كل بُعد تضرب \omterm{def:g6:measure:area}{مساحة} القاعدة في
$2^2 = 4$ والارتفاع في $2$، ومنه \omterm{def:g6:measure:volume}{الحجم} في
$4 \times 2 = 8 = 2^3$. والتكبير بالمعامل $k$ يضرب المساحات في $k^2$
وطولًا آخر في $k$: فتنمو الحجوم بالمعامل $k^3$. و
\omterm{def:g6:lines:objects}{القطعة} العليا في السؤال 12 نسخة مكبّرة من \omterm{def:g8:solids:def}{الهرم} كله
بالمعامل $k = \frac12$، إذن حجمها
$\left(\frac12\right)^3 = \frac18$ من الكلّ --- وهو التفسير
في سطر واحد.
\end{solution}
